\begin{proposition}[6.11.5]
\label{IV.2.s6.11.5-fr}
Sean $A$ un anillo local noetheriano y $M$ un $A$-módulo finitamente generado. Para
todo ideal primo $\mathfrak p$ de $A$, se tiene
\begin{equation}
\operatorname{coprof}_{A_{\mathfrak p}}(M_{\mathfrak p})\leq
\operatorname{coprof}_A(M).
\tag{6.11.5.1}\label{IV.2.s6.11.5.1-fr}
\end{equation}
\end{proposition}

\begin{proof}
Es posible limitarse al caso en que $M_{\mathfrak p}\ne0$. Como $\widehat A$ es
un $A$-módulo fielmente plano (\textbf{0}$_{\rm I}$, 7.3.5), existe un ideal
primo $\mathfrak q$ de $\widehat A$ situado sobre $\mathfrak p$
(\textbf{0}$_{\rm I}$, 6.5.1); como
$\widehat M_{\mathfrak q}=(M\otimes_A\widehat A)\otimes_{\widehat A}\widehat A_{\mathfrak q}
=M\otimes_A\widehat A_{\mathfrak q}$ y
$\widehat A_{\mathfrak q}$ es un $\widehat A$-módulo plano y, por tanto, un
$A$-módulo plano (\textbf{0}$_{\rm I}$, 6.2.1), se sigue de (6.3.2),
aplicado al homomorfismo local
$A_{\mathfrak p}\to\widehat A_{\mathfrak q}$, a $M_{\mathfrak p}$ y a
$\widehat M_{\mathfrak q}=M_{\mathfrak p}\otimes_{A_{\mathfrak p}}\widehat A_{\mathfrak q}$, que se tiene
$\operatorname{coprof}_{A_{\mathfrak p}}(M_{\mathfrak p})\leq
\operatorname{coprof}_{\widehat A_{\mathfrak q}}
(\widehat M_{\mathfrak q})$. Por otra parte,
$X=\operatorname{Spec}(\widehat A)$ es isomorfo a un subesquema de un
esquema regular en virtud del teorema de Cohen (\textbf 0, 19.8.8, (i)). Se deduce,
pues, de (6.11.2) que se tiene
$\operatorname{coprof}_{\widehat A_{\mathfrak q}}
(\widehat M_{\mathfrak q})\leq
\operatorname{coprof}_{\widehat A}(\widehat M)$; finalmente, se sabe que
$\operatorname{coprof}_{\widehat A}(\widehat M)=
\operatorname{coprof}_A(M)$ (\textbf 0, 16.4.10), lo que concluye la demostración de
(6.11.5).
\end{proof}

Esta proposición justifica la definición de la coprofundidad de un
$A$-módulo cuando $A$ no es un anillo local, dada en (5.7.12).
